\documentclass{article}
\usepackage{amssymb}
\usepackage{amsfonts}
\usepackage{amsmath}

\author{Jonathan Kalsky}
\date{April 27 2026}
\title{\normalsize CSCE 222\\ \Large Homework 6}
\linespread{1.5}

\begin{document}
\maketitle
\tableofcontents
\vspace{1cm}

\subsection{6.1}
\subsubsection*{26 - How many strings of four decimal digits}
\subsubsection[26a]{do not contain the same digit twice?}
There are 10 possible digits, and every spot has 1 less possible digit than its previous, or P(10,4).\\
There are 10x9x8x7 = 5040 possible strings according to the product rule to combine these tasks
\subsubsection[26b]{end with an even digit?}
There are 10 possible digits for the first three places, and 5 possible digits for the last place\\
According to the product rule, there are 10x10x10x5 possible ways for these independent tasks to occur\\
Therefore there are 5000 ways to end with an even digit.
\subsubsection[26c]{have exactly three digits that are 9s?}
Like in the stars and bars problem, the non-9 digit is the partition\\
There are 4 spots the unique digit can take place, or C(4,1), and 9 possible digits to fill that spot\\
All other spots are constant (filled with 9's)\\
Therefore there are 4x9, or 36 ways to have exactly three digits that are 9.
\subsubsection*{34 - How many different functions are there from a set with 10 elements to sets with the following numbers of elements?}
\subsubsection[34a]{2 elements}
Each element from the domain has the possibility to map to 2 choices from the codomain. That is the product rule of 2 choices 10 times, or $2^{10}$\\$\therefore 1024$ different functions.
\subsubsection[34b]{3 elements}
Each element from the domain has the possibility to map to 3 choices from the codomain. That is the product rule of 3 choices 10 times, or $3^{10}$\\$\therefore 59049$ different functions.

\subsubsection*{50 - How many bit strings of length seven either begin with
two 0s or end with three 1s?}
\subsubsection{50}
Using the sum rule, we can break this down to 2 problems:\\
\underbar{bit strings of length 7, beggining with two 0's}\\
the two 0's are constant, there is 1 way to fill those two spots, remaining are 5 spots. For each remaining spot there is either a 1 or a 0 (2 choices), so according to the product rule there are $1 \times 2^5$ bit strings.\\
\underbar{bit strings of length 7, ending with three 1's}\\
the three 1's are constant, there is 1 way to fill those three spots, remaining are 4 spots. For each remaining spot, there is either a 1 or a 0 (2 choices), so according to the product rule, there are $1 \times 2^4$ bit strings.\\
\underbar{either}\\
For either of those string formats, we can use the sum rule such that the total number of possible bit strings is $2^5 + 2^4 - 2^2 = 44$. Where $2^2$ (subtraction rule) excludes the double counted strings containing both leading 00 and trailing 111, leaving 2 remaining spots, and 4 possible strings (that are double counted).

\subsection{6.3}
\subsubsection*{20 - How many bit strings of length 10 have}
\subsubsection[20a]{exactly three 0s}
Finding the amount of strings with three 0's and 7 1's, is the same as 10 choose 3 or 7. This is because you only need to find the amount of spots that will be 0's the rest will be (have no choice but being) filled with 1s. C(10,3)=120.
\subsubsection[20b]{more 0s than 1s}
For more 0s than 1s, we would use the sum rule for all strings with more 0 choices than 1: C(10,6) + C(10,7) + C(10,8) + C(10,9) + C(10,10) = 386 strings
\subsubsection[20c]{at least seven 1s}
Similarly, for more than 6 1s chosen, we will use the sum rule: C(10,7) + C(10,8) + C(10,9) + C(10,10) = 176 strings
\subsubsection[20d]{at least three 1s}
Similarly we can add all combinations of 10 choosing more than 2 1s, or subtract less than three 1s from the total string combinations.\\
Total = 2 choices 10 times by the product rule, or $2^{10} = 1024$\\
1024 - C(10,0) - C(10,1) - C(10,2) = 1024-56=968 strings
\subsubsection*{28 - Thirteen people on a softball team show up for a game.}
\subsubsection[28a]{How many ways are there to choose 10 players to take the field?}
Given, 13 and wanting to choose 10 the total combinations = $C(13,10) = \frac{13!}{10!(3)!}=\frac{13\times 12 \times 11}{3\times2} = 286$ ways
\subsubsection[28b]{How many ways are there to assign the 10 positions by selecting players from the 13 people who show up?}
To assign 13 people to 10 spots, there are P(13,10) permutations = $\frac{13!}{3!} = 13\times12\times11\times10\times9\times8\times7\times6\times5\times4=1,037,836,800$ ways
\subsubsection[28c]{Of the 13 people who show up, three are left-handed. How many ways are there to choose 10 players to take the field if at least one of these players must be left-handed?}
From part a we know there are 286 ways to choose 10 players of the 13. To choose at least 1 left handed is the opposite of choosing no lefties (subtraction rule).\\
Thus 286 - C(10,10) = 286-1 = 285 ways.
\subsubsection*{32 - Seven women and nine men are on the faculty in the
mathematics department at a school.}
\subsubsection[32a]{How many ways are there to select a committee of five members of the department if at least one woman must be on the committee?}
Using the subtraction rule:\\
there are a total of C(7+9,5) = 4368 ways to choose 5 members.\\
To choose all 5 as men C(9,5) = 126 ways to not have any woman be on the committee\\
4368-126 = 4242 ways of having at least one woman on board.
\subsubsection[32b]{How many ways are there to select a committee of five members of the department if at least one woman and at least one man must be on the committee?}
For the committee to also have at least one man, we need to continue the subtraction rule from part b.\\
4242 - C(7,5) discounts all the cases of a committee full of women. Since there are no duplicates between a committee of all men and all women, we dont need to subtract duplicate ways.\\
Total = 4242 - 21 = 4221 ways.



\end{document}